%! Factoring: GCF and Grouping — Unit 1, Lesson 1.3 — Guided Notes
\documentclass[10pt]{article}
\usepackage{atda-article}
\usepackage{atda-boxes}
\newcommand{\TallMath}[1]{$\displaystyle #1\rule[-1.4em]{0pt}{3.2em}$}
\setlength{\workrowsep}{6pt}

\begin{document}

\pageheader{Unit 1, Lesson 1.3}{Guided Notes}
% NAMESTRIP: no \namedateperiod here. The student writes their name once, on the
% cover this component is stapled behind. See references/conventions.md.

\vspace{0.15in}

\begin{objectivebox}
\small
By the end of this lesson, I will be able to\dots
\begin{enumerate}[leftmargin=1.5em, itemsep=2pt, topsep=3pt]
  \item find the \textbf{greatest common factor} of a polynomial in one or two variables and pull
        it out front;
  \item factor a \textbf{four-term} polynomial \textbf{by grouping}, including when a pair needs a
        negative factored out;
  \item factor \textbf{completely} --- GCF first, then group --- and check by multiplying back;
  \item read a factored area or volume model to recover the \textbf{dimensions} of the figure it
        describes, and say what those dimensions mean.
\end{enumerate}
\end{objectivebox}

\vspace{0.10in}

\boxguard
\begin{vocabbox}
\small
Fill in each term as we build it together.
\termblanklong{Factoring}
\termblanklong{Greatest common factor (GCF)}
\termblanklong{Factored completely}
\termblanklong{Factoring by grouping}
\termblanklong{Common binomial factor}
\end{vocabbox}

\vspace{0.10in}

\boxguard[12]
\begin{hookbox}
\small
\textbf{The blueprint with the dimensions erased.} A contractor's plan gives the area of a
rectangular patio as
\[ A(x) = 24x^2+36x \quad \text{square feet,} \]
where $x$ is the number of feet the owner chooses to widen the design. The tile crew does not need
the area --- they need the two \textbf{side lengths}, and multiplying erased them.

Can you get the dimensions back out of the area? Decide first, then check it.

\writelines{2}
\end{hookbox}

\vspace{0.10in}

\boxguard[24]
\begin{notesbox}{1. The Greatest Common Factor --- Pull Out the Largest Shared Piece}
\small
\textbf{Factoring is multiplying, run backwards.} The distributive property built
$6(2x+5)$ into $12x+30$; factoring takes $12x+30$ back to $6(2x+5)$.

To find the GCF of a polynomial, take the GCF of the two pieces separately:
\begin{itemize}[leftmargin=1.3em, itemsep=2pt, topsep=3pt]
  \item \textbf{Numbers:} the largest integer that divides every coefficient.
  \item \textbf{Variables:} each variable that appears in \emph{every} term, carrying its
        \emph{smallest} exponent.
\end{itemize}

\renewcommand{\arraystretch}{1.9}
\begin{tabularx}{\linewidth}{@{}p{4.0cm} p{2.6cm} X@{}}
\rowcolor{mist}
\textbf{Polynomial} & \textbf{GCF} & \textbf{Factored form} \\
\hline
$8x+20$                  & \blank{1.8cm} & \blank{3.4cm} \\
$15x^3-25x^2$            & \blank{1.8cm} & \blank{3.4cm} \\
$6a^2b+9ab^2$            & \blank{1.8cm} & \blank{3.4cm} \\
$14x^3-7x^2+21x$         & \blank{1.8cm} & \blank{3.4cm} \\
$-4y^3-12y^2$            & \blank{1.8cm} & \blank{3.4cm} \\
\end{tabularx}

\vspace{6pt}
\textbf{Careful --- this is Lesson 1.2's warning again.} $\sqrt{72}=\sqrt4\cdot\sqrt{18}$ was true
but not simplified, because a perfect square was still hiding inside. In the same way
$24x^2+36x=6x(4x+6)$ is a true statement and a wrong answer: the leftover $4x+6$ still shares a
factor of $2$. A factorization is \blank{2.6cm} only when the terms left inside share nothing but
$1$.

\textbf{Back to the hook} --- the patio, and the check that costs ten seconds:
\begin{work}
  24x^2+36x &= 12x(2x) + 12x(3) \\
            &= 12x(2x+3) \\
  \text{Check: } 12x(2x+3) &= 24x^2+36x \quad\checkmark
\end{work}
The crew reads off a width of $12x$ feet and a length of $(2x+3)$ feet.
\end{notesbox}

\vspace{0.10in}

\boxguard[22]
\begin{notesbox}{2. Four Terms, Nothing Shared by All of Them --- Group}
\small
Sometimes no factor is common to \emph{every} term. Warm-up item 4 showed the way out: split the
four terms into two pairs, pull the GCF out of each pair, and look at what is left.
\begin{work}
  2x^3+6x^2+5x+15 &= \left(2x^3+6x^2\right) + (5x+15) \\
                  &= 2x^2(x+3) + 5(x+3) \\
                  &= (x+3)\left(2x^2+5\right)
\end{work}
The pairs each left behind $(x+3)$ --- a \textbf{common binomial factor} --- and it comes out front
exactly the way a monomial does. \textbf{The two leftovers must be identical}; if they are not, the
\blank{2.4cm} was wrong, so try the other one (or write the terms in descending order first).

\renewcommand{\arraystretch}{1.9}
\begin{tabularx}{\linewidth}{@{}p{4.3cm} X | X@{}}
\rowcolor{mist}
\textbf{Four-term polynomial} & \textbf{GCF of each pair} & \textbf{Factored form} \\
\hline
$x^3+4x^2+3x+12$    & \blank{3.2cm} & \blank{3.0cm} \\
$6x^3-9x^2+4x-6$    & \blank{3.2cm} & \blank{3.0cm} \\
$x^3-5x^2-3x+15$    & \blank{3.2cm} & \blank{3.0cm} \\
\end{tabularx}

\vspace{6pt}
\textbf{Watch the sign on the third one.} When the third term is negative, factor out a
\emph{negative} GCF from the second pair --- $-3x+15 = -3(x-5)$, not $-3(x+5)$. Getting this
backwards is what makes two leftovers disagree that should have matched.
\end{notesbox}

\vspace{0.10in}

\boxguard[20]
\begin{notesbox}{3. Factoring Completely --- The Order to Work In}
\small
``Factor completely'' means keep going until nothing else can come out. Work in this order, every
time:
\begin{enumerate}[leftmargin=1.4em, itemsep=2pt, topsep=3pt]
  \item \textbf{GCF first, always.} It makes every number after it smaller.
  \item \textbf{Count the terms.} Four terms $\Rightarrow$ group. (Three and two terms get their
        own methods in Lessons 1.4 and 1.5.)
  \item \textbf{Check every factor} for something still shared, then multiply back to confirm.
\end{enumerate}
\begin{work}
  2x^3+10x^2+6x+30 &= 2\left(x^3+5x^2+3x+15\right) \\
                   &= 2\left[x^2(x+5) + 3(x+5)\right] \\
                   &= 2(x+5)\left(x^2+3\right)
\end{work}
Skipping the GCF here does not make the problem impossible --- it makes it uglier. Grouping
$2x^3+10x^2+6x+30$ directly gives $2x^2(x+5)+6(x+5)$, and the $2$ still has to come out at the end.
\end{notesbox}

\vspace{0.10in}

\boxguard[20]
\begin{practicebox}
\small
\textbf{The deck job.} The same contractor quotes a rectangular deck with area
$A(x)=30x^2+42x$ square feet. The plans say one side measures $6x$ feet.

\begin{enumerate}[leftmargin=1.6em, itemsep=6pt, topsep=4pt]
  \item Factor $A(x)$ completely, and use it to find the other side of the deck.
\begin{work}
  30x^2+42x &= 6x(5x) + 6x(7) \\
            &= 6x(5x+7)
\end{work}

  \item The owner picks $x=4$. Give both side lengths and the area, in feet, with units.

\writeline

  \item A second deck is built as two rectangular sections, with total area
        $A(x)=5x^3+15x^2+2x+6$. Factor it, then explain what the shared factor tells you about
        how the two sections sit next to each other.
\begin{work}
  5x^3+15x^2+2x+6 &= 5x^2(x+3) + 2(x+3) \\
                  &= (x+3)\left(5x^2+2\right)
\end{work}
\writelines{2}
\end{enumerate}
\end{practicebox}

\end{document}
